% Options for packages loaded elsewhere
\PassOptionsToPackage{unicode}{hyperref}
\PassOptionsToPackage{hyphens}{url}
\documentclass[
]{article}
\usepackage{xcolor}
\usepackage[margin=1in]{geometry}
\usepackage{amsmath,amssymb}
\setcounter{secnumdepth}{-\maxdimen} % remove section numbering
\usepackage{iftex}
\ifPDFTeX
  \usepackage[T1]{fontenc}
  \usepackage[utf8]{inputenc}
  \usepackage{textcomp} % provide euro and other symbols
\else % if luatex or xetex
  \usepackage{unicode-math} % this also loads fontspec
  \defaultfontfeatures{Scale=MatchLowercase}
  \defaultfontfeatures[\rmfamily]{Ligatures=TeX,Scale=1}
\fi
\usepackage{lmodern}
\ifPDFTeX\else
  % xetex/luatex font selection
\fi
% Use upquote if available, for straight quotes in verbatim environments
\IfFileExists{upquote.sty}{\usepackage{upquote}}{}
\IfFileExists{microtype.sty}{% use microtype if available
  \usepackage[]{microtype}
  \UseMicrotypeSet[protrusion]{basicmath} % disable protrusion for tt fonts
}{}
\makeatletter
\@ifundefined{KOMAClassName}{% if non-KOMA class
  \IfFileExists{parskip.sty}{%
    \usepackage{parskip}
  }{% else
    \setlength{\parindent}{0pt}
    \setlength{\parskip}{6pt plus 2pt minus 1pt}}
}{% if KOMA class
  \KOMAoptions{parskip=half}}
\makeatother
\usepackage{color}
\usepackage{fancyvrb}
\newcommand{\VerbBar}{|}
\newcommand{\VERB}{\Verb[commandchars=\\\{\}]}
\DefineVerbatimEnvironment{Highlighting}{Verbatim}{commandchars=\\\{\}}
% Add ',fontsize=\small' for more characters per line
\usepackage{framed}
\definecolor{shadecolor}{RGB}{248,248,248}
\newenvironment{Shaded}{\begin{snugshade}}{\end{snugshade}}
\newcommand{\AlertTok}[1]{\textcolor[rgb]{0.94,0.16,0.16}{#1}}
\newcommand{\AnnotationTok}[1]{\textcolor[rgb]{0.56,0.35,0.01}{\textbf{\textit{#1}}}}
\newcommand{\AttributeTok}[1]{\textcolor[rgb]{0.13,0.29,0.53}{#1}}
\newcommand{\BaseNTok}[1]{\textcolor[rgb]{0.00,0.00,0.81}{#1}}
\newcommand{\BuiltInTok}[1]{#1}
\newcommand{\CharTok}[1]{\textcolor[rgb]{0.31,0.60,0.02}{#1}}
\newcommand{\CommentTok}[1]{\textcolor[rgb]{0.56,0.35,0.01}{\textit{#1}}}
\newcommand{\CommentVarTok}[1]{\textcolor[rgb]{0.56,0.35,0.01}{\textbf{\textit{#1}}}}
\newcommand{\ConstantTok}[1]{\textcolor[rgb]{0.56,0.35,0.01}{#1}}
\newcommand{\ControlFlowTok}[1]{\textcolor[rgb]{0.13,0.29,0.53}{\textbf{#1}}}
\newcommand{\DataTypeTok}[1]{\textcolor[rgb]{0.13,0.29,0.53}{#1}}
\newcommand{\DecValTok}[1]{\textcolor[rgb]{0.00,0.00,0.81}{#1}}
\newcommand{\DocumentationTok}[1]{\textcolor[rgb]{0.56,0.35,0.01}{\textbf{\textit{#1}}}}
\newcommand{\ErrorTok}[1]{\textcolor[rgb]{0.64,0.00,0.00}{\textbf{#1}}}
\newcommand{\ExtensionTok}[1]{#1}
\newcommand{\FloatTok}[1]{\textcolor[rgb]{0.00,0.00,0.81}{#1}}
\newcommand{\FunctionTok}[1]{\textcolor[rgb]{0.13,0.29,0.53}{\textbf{#1}}}
\newcommand{\ImportTok}[1]{#1}
\newcommand{\InformationTok}[1]{\textcolor[rgb]{0.56,0.35,0.01}{\textbf{\textit{#1}}}}
\newcommand{\KeywordTok}[1]{\textcolor[rgb]{0.13,0.29,0.53}{\textbf{#1}}}
\newcommand{\NormalTok}[1]{#1}
\newcommand{\OperatorTok}[1]{\textcolor[rgb]{0.81,0.36,0.00}{\textbf{#1}}}
\newcommand{\OtherTok}[1]{\textcolor[rgb]{0.56,0.35,0.01}{#1}}
\newcommand{\PreprocessorTok}[1]{\textcolor[rgb]{0.56,0.35,0.01}{\textit{#1}}}
\newcommand{\RegionMarkerTok}[1]{#1}
\newcommand{\SpecialCharTok}[1]{\textcolor[rgb]{0.81,0.36,0.00}{\textbf{#1}}}
\newcommand{\SpecialStringTok}[1]{\textcolor[rgb]{0.31,0.60,0.02}{#1}}
\newcommand{\StringTok}[1]{\textcolor[rgb]{0.31,0.60,0.02}{#1}}
\newcommand{\VariableTok}[1]{\textcolor[rgb]{0.00,0.00,0.00}{#1}}
\newcommand{\VerbatimStringTok}[1]{\textcolor[rgb]{0.31,0.60,0.02}{#1}}
\newcommand{\WarningTok}[1]{\textcolor[rgb]{0.56,0.35,0.01}{\textbf{\textit{#1}}}}
\usepackage{longtable,booktabs,array}
\usepackage{calc} % for calculating minipage widths
% Correct order of tables after \paragraph or \subparagraph
\usepackage{etoolbox}
\makeatletter
\patchcmd\longtable{\par}{\if@noskipsec\mbox{}\fi\par}{}{}
\makeatother
% Allow footnotes in longtable head/foot
\IfFileExists{footnotehyper.sty}{\usepackage{footnotehyper}}{\usepackage{footnote}}
\makesavenoteenv{longtable}
\usepackage{graphicx}
\makeatletter
\newsavebox\pandoc@box
\newcommand*\pandocbounded[1]{% scales image to fit in text height/width
  \sbox\pandoc@box{#1}%
  \Gscale@div\@tempa{\textheight}{\dimexpr\ht\pandoc@box+\dp\pandoc@box\relax}%
  \Gscale@div\@tempb{\linewidth}{\wd\pandoc@box}%
  \ifdim\@tempb\p@<\@tempa\p@\let\@tempa\@tempb\fi% select the smaller of both
  \ifdim\@tempa\p@<\p@\scalebox{\@tempa}{\usebox\pandoc@box}%
  \else\usebox{\pandoc@box}%
  \fi%
}
% Set default figure placement to htbp
\def\fps@figure{htbp}
\makeatother
\setlength{\emergencystretch}{3em} % prevent overfull lines
\providecommand{\tightlist}{%
  \setlength{\itemsep}{0pt}\setlength{\parskip}{0pt}}
\usepackage{bookmark}
\IfFileExists{xurl.sty}{\usepackage{xurl}}{} % add URL line breaks if available
\urlstyle{same}
\hypersetup{
  hidelinks,
  pdfcreator={LaTeX via pandoc}}

\author{}
\date{\vspace{-2.5em}}

\begin{document}

\section{📚 RepositorioMatematicasICFES\_R\_Exams (Filosofía
2025)}\label{repositoriomatematicasicfes_r_exams-filosofuxeda-2025}

\href{https://github.com/alvaretto/proyecto-r-exams-icfes-matematicas-optimizado}{\pandocbounded{\includegraphics[keepaspectratio]{https://img.shields.io/badge/Estado-Activo-brightgreen}}}
\href{https://www.r-exams.org/}{\pandocbounded{\includegraphics[keepaspectratio]{https://img.shields.io/badge/R--exams-Compatible-orange}}}
\href{https://github.com/alvaretto/proyecto-r-exams-icfes-matematicas-optimizado}{\pandocbounded{\includegraphics[keepaspectratio]{https://img.shields.io/badge/Calidad-ICFES_2025-success}}}
\href{https://github.com/alvaretto/proyecto-r-exams-icfes-matematicas-optimizado}{\pandocbounded{\includegraphics[keepaspectratio]{https://img.shields.io/badge/Metodología-Sistema_Condicional-blue}}}

\textbf{Sistema integrado para la generación de ejercicios matemáticos
tipo ICFES en formato R-exams a partir de imágenes, siguiendo la
filosofía ``Matemáticas ICFES 2025''.}

Este repositorio implementa un flujo de trabajo completo y automatizado
que transforma imágenes de problemas matemáticos en ejercicios
\texttt{.Rmd} interactivos, parametrizados y de alta calidad, listos
para ser compilados con R-exams.

\begin{center}\rule{0.5\linewidth}{0.5pt}\end{center}

\subsection{\texorpdfstring{🎯 \textbf{Filosofía del
Sistema}}{🎯 Filosofía del Sistema}}\label{filosofuxeda-del-sistema}

El proyecto se basa en un flujo de trabajo de 3 fases diseñado para
maximizar la calidad y la eficiencia:

\begin{enumerate}
\def\labelenumi{\arabic{enumi}.}
\tightlist
\item
  \textbf{📥 Fase 1: Entrada de Datos y Análisis Automático}

  \begin{itemize}
  \tightlist
  \item
    Recepción de una imagen (\texttt{.png}) con un problema matemático.
  \item
    \textbf{Sistema Condicional Automático}: Detección inteligente de
    contenido gráfico para activar el flujo de trabajo adecuado.
  \item
    Protocolo de conversión para imágenes que no siguen el formato ICFES
    estándar.
  \end{itemize}
\item
  \textbf{⚙️ Fase 2: Procesamiento y Generación}

  \begin{itemize}
  \tightlist
  \item
    \textbf{Flujo A (Sin Gráficas)}: Proceso estándar de 8 fases para la
    generación del \texttt{.Rmd}.
  \item
    \textbf{Flujo B (Con Gráficas)}: Activación del
    \textbf{Agente-Graficador Especializado TikZ} para replicar la
    imagen con una fidelidad visual superior al 98\%.
  \item
    Generación de código \texttt{.Rmd} completo, siguiendo una
    estructura obligatoria y validada.
  \end{itemize}
\item
  \textbf{🔄 Fase 3: Iteración y Mejora Continua}

  \begin{itemize}
  \tightlist
  \item
    Captura de retroalimentación para la corrección de errores.
  \item
    Implementación de mejoras basadas en patrones identificados para
    optimizar la generación de código.
  \end{itemize}
\end{enumerate}

\begin{center}\rule{0.5\linewidth}{0.5pt}\end{center}

\subsection{\texorpdfstring{✨ \textbf{Características
Principales}}{✨ Características Principales}}\label{caracteruxedsticas-principales}

\begin{itemize}
\tightlist
\item
  \textbf{🤖 Sistema Condicional Automático}: Activa flujos de trabajo
  especializados según el contenido de la imagen.
\item
  \textbf{🎨 Agente-Graficador TikZ}: Replica imágenes con una fidelidad
  visual del 98\%+, manejando colores RGB precisos, posicionamiento
  sistemático y características avanzadas.
\item
  \textbf{🧩 Soporte Híbrido}: Genera ejercicios que integran R, Python
  (vía \texttt{reticulate}), y LaTeX/TikZ para visualizaciones
  complejas.
\item
  \textbf{✅ Calidad ICFES Garantizada}: Cumple con metadatos
  obligatorios, sistema avanzado de distractores y criterios de calidad
  rigurosos.
\item
  \textbf{🔧 Protocolo Anti-Errores}: Utiliza una base de ``ejemplos
  funcionales'' para prevenir errores de implementación, asegurando que
  todo el código generado sea robusto y compilable.
\item
  \textbf{✔️ Pruebas de Diversidad}: Generación mínima de 300 versiones
  únicas por ejercicio, verificada con \texttt{testthat}.
\end{itemize}

\begin{center}\rule{0.5\linewidth}{0.5pt}\end{center}

\subsection{\texorpdfstring{🛠️ \textbf{Tecnologías
Utilizadas}}{🛠️ Tecnologías Utilizadas}}\label{tecnologuxedas-utilizadas}

\begin{itemize}
\tightlist
\item
  \textbf{Motor Principal}: \textbf{R} (≥ 4.0) y \textbf{R-exams}.
\item
  \textbf{Visualización}:

  \begin{itemize}
  \tightlist
  \item
    \textbf{LaTeX/TikZ}: Para diagramas, figuras geométricas y
    notaciones matemáticas de alta calidad.
  \item
    \textbf{Python (\texttt{matplotlib}, \texttt{numpy})}: Integrado con
    \texttt{reticulate} para gráficos avanzados.
  \item
    \textbf{R (\texttt{ggplot2})}: Para análisis y gráficos
    estadísticos.
  \end{itemize}
\item
  \textbf{Entorno}: VSCode con herramientas de IA, sobre Manjaro Plasma
  KDE.
\item
  \textbf{Control de Versiones}: Git y GitHub.
\end{itemize}

\begin{center}\rule{0.5\linewidth}{0.5pt}\end{center}

\subsection{\texorpdfstring{📁 \textbf{Estructura del
Repositorio}}{📁 Estructura del Repositorio}}\label{estructura-del-repositorio}

\begin{verbatim}
RepositorioMatematicasICFES_R_Exams/
├── 📂 01-Algebra-Y-Calculo/
├── 📂 05-Geometria/
├── 📂 06-Estadistica-Y-Probabilidad/
├── 🏆 A-Produccion/                        # ⭐ Archivos listos para producción
│   ├── 📂 Ejemplos-Funcionales-Rmd/       # (OBLIGATORIO) Base de código validado
│   ├── 📂 01-Numeros-Reales/              # Ejercicios de producción por tema
│   ├── 📂 02-Funciones/
│   ├── 📂 05-Geometría/
│   └── 📂 06-Estadística-Y-Probabilidad/
├── 🛠️ Auxiliares/
│   ├── 📂 Agente-Graficador-TikZ/         # Laboratorio del agente TikZ
│   ├── 📂 Estrategia-Avanzada-de-Replicas/ # Templates TikZ profesionales
│   └── 📂 Python-Documentation/             # Guías para integración con Python
├── 🧪 Lab-Manjaro/                         # Ejercicios en desarrollo y pruebas
├── 📄 .gitignore
├── 📖 README.md                            # Este archivo
└── 📖 walkthrough.md                       # Guía de uso detallada
\end{verbatim}

\begin{center}\rule{0.5\linewidth}{0.5pt}\end{center}

\subsection{\texorpdfstring{🏆 \textbf{Carpeta A-Produccion: Archivos
Listos para
Uso}}{🏆 Carpeta A-Produccion: Archivos Listos para Uso}}\label{carpeta-a-produccion-archivos-listos-para-uso}

\subsubsection{\texorpdfstring{📋
\textbf{Descripción}}{📋 Descripción}}\label{descripciuxf3n}

La carpeta \textbf{\texttt{A-Produccion/}} contiene archivos
\texttt{.Rmd} y subproyectos funcionales que han sido validados y están
\textbf{listos para producción}. Estos archivos representan el estándar
de calidad del proyecto y pueden utilizarse directamente o con ajustes
mínimos.

\subsubsection{\texorpdfstring{✨ \textbf{Características de los
Archivos en
A-Produccion}}{✨ Características de los Archivos en A-Produccion}}\label{caracteruxedsticas-de-los-archivos-en-a-produccion}

\begin{itemize}
\tightlist
\item
  ✅ \textbf{Completamente funcionales}: Compilables sin errores en
  R-exams
\item
  ✅ \textbf{Validados}: Han pasado todas las pruebas de calidad y
  diversidad (300+ versiones)
\item
  ✅ \textbf{Documentados}: Incluyen comentarios y estructura clara
\item
  ✅ \textbf{Optimizados}: Siguen las mejores prácticas del proyecto
\item
  ✅ \textbf{Modelos de referencia}: Pueden usarse como plantillas para
  nuevos ejercicios
\end{itemize}

\subsubsection{\texorpdfstring{📂 \textbf{Contenido de
A-Produccion}}{📂 Contenido de A-Produccion}}\label{contenido-de-a-produccion}

\begin{verbatim}
A-Produccion/
├── 📂 Ejemplos-Funcionales-Rmd/       # ⭐ Base de código validado (CONSULTA OBLIGATORIA)
│   ├── Plantillas/                    # Templates por tipo de ejercicio
│   │   ├── Rmd/                       # Plantillas .Rmd (schoice, cloze, etc.)
│   │   ├── Rnw/                       # Plantillas .Rnw
│   │   └── TikZ-Documentation/        # Documentación y ejemplos TikZ
│   └── [Ejercicios validados]         # Ejercicios funcionales de referencia
├── 📂 01-Numeros-Reales/              # Ejercicios de producción por tema
├── 📂 02-Funciones/
├── 📂 05-Geometría/
└── 📂 06-Estadística-Y-Probabilidad/
\end{verbatim}

\subsubsection{\texorpdfstring{🎯 \textbf{Uso
Recomendado}}{🎯 Uso Recomendado}}\label{uso-recomendado}

\begin{enumerate}
\def\labelenumi{\arabic{enumi}.}
\tightlist
\item
  \textbf{Como referencia}: Consultar antes de crear nuevos ejercicios
\item
  \textbf{Como plantilla}: Copiar y adaptar para nuevos problemas
  similares
\item
  \textbf{Como validación}: Comparar tu código con estos ejemplos
  funcionales
\item
  \textbf{Como aprendizaje}: Estudiar las mejores prácticas
  implementadas
\end{enumerate}

\subsubsection{\texorpdfstring{⚠️
\textbf{Importante}}{⚠️ Importante}}\label{importante}

\begin{quote}
\textbf{Regla de Oro}: ``Si no está en los ejemplos funcionales, no
improvises.''

Antes de escribir cualquier código nuevo, es \textbf{OBLIGATORIO}
consultar los ejercicios en
\texttt{/A-Produccion/Ejemplos-Funcionales-Rmd/} para seguir los
patrones validados.
\end{quote}

\begin{center}\rule{0.5\linewidth}{0.5pt}\end{center}

\subsection{\texorpdfstring{📦 \textbf{Cambio de Ubicación:
Ejemplos-Funcionales-Rmd}}{📦 Cambio de Ubicación: Ejemplos-Funcionales-Rmd}}\label{cambio-de-ubicaciuxf3n-ejemplos-funcionales-rmd}

\subsubsection{\texorpdfstring{🔄 \textbf{Reorganización del
Repositorio}}{🔄 Reorganización del Repositorio}}\label{reorganizaciuxf3n-del-repositorio}

La carpeta \textbf{\texttt{Ejemplos-Funcionales-Rmd}} ha sido
\textbf{reubicada} para mejorar la organización del proyecto y destacar
su importancia como base de código de producción.

\subsubsection{\texorpdfstring{📍
\textbf{Ubicaciones}}{📍 Ubicaciones}}\label{ubicaciones}

\begin{longtable}[]{@{}
  >{\raggedright\arraybackslash}p{(\linewidth - 2\tabcolsep) * \real{0.5000}}
  >{\raggedright\arraybackslash}p{(\linewidth - 2\tabcolsep) * \real{0.5000}}@{}}
\toprule\noalign{}
\begin{minipage}[b]{\linewidth}\raggedright
Aspecto
\end{minipage} & \begin{minipage}[b]{\linewidth}\raggedright
Detalle
\end{minipage} \\
\midrule\noalign{}
\endhead
\bottomrule\noalign{}
\endlastfoot
\textbf{Ubicación Anterior} &
\texttt{/Auxiliares/Ejemplos-Funcionales-Rmd/} \\
\textbf{Ubicación Nueva} &
\texttt{/A-Produccion/Ejemplos-Funcionales-Rmd/} \\
\textbf{Fecha del Cambio} & Diciembre 2025 \\
\textbf{Razón del Cambio} & Destacar archivos de producción y mejorar
organización \\
\end{longtable}

\subsubsection{\texorpdfstring{🎯 \textbf{Propósito del
Cambio}}{🎯 Propósito del Cambio}}\label{propuxf3sito-del-cambio}

\begin{enumerate}
\def\labelenumi{\arabic{enumi}.}
\tightlist
\item
  \textbf{Mayor visibilidad}: Los ejemplos funcionales ahora están en
  una carpeta de producción destacada
\item
  \textbf{Mejor organización}: Separación clara entre archivos de
  producción y auxiliares
\item
  \textbf{Facilitar acceso}: Ubicación más intuitiva para archivos de
  referencia obligatoria
\item
  \textbf{Estandarización}: Agrupar todos los archivos listos para
  producción en un solo lugar
\end{enumerate}

\subsubsection{\texorpdfstring{📝 \textbf{Actualización de
Referencias}}{📝 Actualización de Referencias}}\label{actualizaciuxf3n-de-referencias}

Todas las referencias en el código, documentación y configuraciones han
sido actualizadas para reflejar la nueva ubicación:

\begin{itemize}
\tightlist
\item
  ✅ Archivos de configuración (\texttt{.json}, \texttt{.yaml})
\item
  ✅ Documentación (\texttt{.md})
\item
  ✅ Reglas de IA (\texttt{.cursor/rules/}, \texttt{.augment/rules/})
\item
  ✅ Scripts y herramientas
\item
  ✅ README y guías de uso
\end{itemize}

\subsubsection{\texorpdfstring{🔧 \textbf{Acción
Requerida}}{🔧 Acción Requerida}}\label{acciuxf3n-requerida}

Si tienes scripts o configuraciones personales que referencian la
ubicación antigua, actualízalas a:

\begin{verbatim}
/A-Produccion/Ejemplos-Funcionales-Rmd/
\end{verbatim}

\begin{center}\rule{0.5\linewidth}{0.5pt}\end{center}

\subsection{\texorpdfstring{🚀 \textbf{Instalación y
Uso}}{🚀 Instalación y Uso}}\label{instalaciuxf3n-y-uso}

\subsubsection{\texorpdfstring{\textbf{Requisitos
Previos}}{Requisitos Previos}}\label{requisitos-previos}

Asegúrate de tener instalado R (≥ 4.0), una distribución de LaTeX (se
recomienda \texttt{tinytex}) y Python.

\begin{Shaded}
\begin{Highlighting}[]
\CommentTok{\# 1. Instalar paquetes esenciales de R}
\FunctionTok{install.packages}\NormalTok{(}\FunctionTok{c}\NormalTok{(}\StringTok{"exams"}\NormalTok{, }\StringTok{"tidyverse"}\NormalTok{, }\StringTok{"ggplot2"}\NormalTok{, }\StringTok{"knitr"}\NormalTok{, }\StringTok{"reticulate"}\NormalTok{, }\StringTok{"testthat"}\NormalTok{, }\StringTok{"tinytex"}\NormalTok{))}

\CommentTok{\# 2. Configurar TinyTeX}
\NormalTok{tinytex}\SpecialCharTok{::}\FunctionTok{install\_tinytex}\NormalTok{()}

\CommentTok{\# 3. Configurar Python para reticulate}
\FunctionTok{library}\NormalTok{(reticulate)}
\FunctionTok{use\_python}\NormalTok{(}\StringTok{"/usr/bin/python3"}\NormalTok{, }\AttributeTok{required =} \ConstantTok{TRUE}\NormalTok{) }\CommentTok{\# Ajusta la ruta a tu ejecutable de Python}
\end{Highlighting}
\end{Shaded}

\subsubsection{\texorpdfstring{\textbf{Uso del
Sistema}}{Uso del Sistema}}\label{uso-del-sistema}

El sistema está diseñado para ser operado mediante \textbf{Skills de
Claude Code} que automatizan el workflow completo.

\paragraph{\texorpdfstring{\textbf{Skills Disponibles} (Workflow
Automatizado)}{Skills Disponibles (Workflow Automatizado)}}\label{skills-disponibles-workflow-automatizado}

El proyecto incluye skills configurados en \texttt{.claude/skills/} para
automatizar cada fase del workflow:

\textbf{Workflow Principal:}

\begin{itemize}
\tightlist
\item
  \texttt{/analizar-icfes} - Análisis ICFES de imagen según 6
  dimensiones (Fase 1)
\item
  \texttt{/generar-schoice} - Generar ejercicio de selección única con
  nomenclatura obligatoria y carpeta estructurada
\item
  \texttt{/generar-cloze} - Generar ejercicio tipo CLOZE con
  configuración de tolerancias apropiadas
\item
  \texttt{/promover-ejercicio} - Promoción a carpeta de producción (Fase
  7)
\end{itemize}

\textbf{Skills de Soporte:}

\begin{itemize}
\tightlist
\item
  \texttt{/corregir-error-imagen} - Corrección automática de errores
  TikZ
\item
  \texttt{/validar-diversidad} - Validar 300+ versiones únicas
\item
  \texttt{/validar-icfes} - Validar metadatos y estructura R-exams
\end{itemize}

\textbf{Skills del Graficador Experto v2.0:}

\begin{itemize}
\tightlist
\item
  \texttt{/analizar-imagen} - Análisis visual detallado con estado
  persistente
\item
  \texttt{/generar-tikz} - Generación TikZ/LaTeX con métricas
  cuantitativas
\item
  \texttt{/generar-python} - Generación Python/Matplotlib con
  transferencia de conocimiento
\item
  \texttt{/generar-r} - Generación R/ggplot2 con lecciones aprendidas
\item
  \texttt{/comparar} - Comparación visual con puntuación 0-100 puntos
\item
  \texttt{/iterar} - Refinamiento iterativo con contador de iteraciones
\item
  \texttt{/exportar} - Exportación completa con estadísticas y carpeta
  nomenclada
\item
  \texttt{/estado} - Visualización de progreso del workflow en tiempo
  real
\item
  \texttt{/auto-iterar} - Iteración automática hasta umbral de similitud
\end{itemize}

\paragraph{\texorpdfstring{\textbf{Ejemplo de Uso con
Skills}}{Ejemplo de Uso con Skills}}\label{ejemplo-de-uso-con-skills}

\begin{Shaded}
\begin{Highlighting}[]
\CommentTok{\# 1. Analizar imagen de ejercicio ICFES}
\ExtensionTok{/analizar{-}icfes}\NormalTok{ imagen\_ejercicio.png}

\CommentTok{\# 2. Si tiene gráficos complejos, usar Graficador Experto}
\ExtensionTok{/analizar{-}imagen}\NormalTok{ grafico\_matematico.png}
\ExtensionTok{/generar{-}r}  \CommentTok{\# Genera código R (recomendado para R{-}exams)}
\ExtensionTok{/comparar}   \CommentTok{\# Obtiene puntuación 0{-}100}
\ExtensionTok{/auto{-}iterar}\NormalTok{ r 95 10  }\CommentTok{\# Itera hasta 95+ puntos (máx 10 iteraciones)}

\CommentTok{\# 3. Generar ejercicio SCHOICE (pregunta al usuario qué versión gráfica usar)}
\ExtensionTok{/generar{-}schoice}

\CommentTok{\# 4. Validar diversidad de versiones}
\ExtensionTok{/validar{-}diversidad}\NormalTok{ archivo\_generado.Rmd}

\CommentTok{\# 5. Promover a producción después de validar}
\ExtensionTok{/promover{-}ejercicio}\NormalTok{ archivo\_generado.Rmd}
\end{Highlighting}
\end{Shaded}

\paragraph{\texorpdfstring{\textbf{Comandos Tradicionales
(Alternativos)}}{Comandos Tradicionales (Alternativos)}}\label{comandos-tradicionales-alternativos}

También puedes usar comandos en lenguaje natural:

\begin{quote}
``Aplica el sistema condicional automático a esta imagen PNG para
detectar contenido gráfico y activar el flujo apropiado.''
\end{quote}

Para problemas que requieren replicación gráfica directa:

\begin{quote}
``Activa el Agente-Graficador Especializado TikZ para replicar esta
imagen con 98\%+ fidelidad visual.''
\end{quote}

El sistema se encargará de analizar la imagen, seleccionar el flujo
correcto y generar el archivo \texttt{.Rmd} correspondiente en el
directorio de trabajo.

\begin{center}\rule{0.5\linewidth}{0.5pt}\end{center}

\subsection{\texorpdfstring{🤖 \textbf{Configuración de Claude
Code}}{🤖 Configuración de Claude Code}}\label{configuraciuxf3n-de-claude-code}

El proyecto incluye configuración completa en \texttt{.claude/} para
automatizar el workflow:

\begin{verbatim}
.claude/
├── CLAUDE.md              # Memory file principal del proyecto (v2.1)
├── rules/                 # Reglas modulares (según doc oficial nov 2025)
│   ├── ciclo-validacion.md           # Ciclo de Validación Automática (OBLIGATORIO)
│   ├── codigo-rmd.md                 # Reglas para código R/Markdown
│   └── documentacion-verificada.md   # Principio de Documentación Verificada
├── settings.json          # Hooks y configuración global
├── settings.local.json    # Permisos para skills
├── schemas/               # Esquemas JSON para estado persistente
│   ├── workflow_state.schema.json
│   ├── analisis_inicial.schema.json
│   ├── metricas_similitud.schema.json
│   └── lecciones_aprendidas.schema.json
├── skills/                # Skills del workflow automatizado
│   ├── analizar-icfes/
│   ├── generar-schoice/
│   ├── generar-cloze/
│   ├── promover-ejercicio/
│   ├── corregir-error-imagen/
│   ├── validar-diversidad/
│   ├── validar-icfes/
│   └── [Graficador Experto v2.0]
│       ├── analizar-imagen-matematica/
│       ├── generar-tikz/
│       ├── generar-python/
│       ├── generar-r/
│       ├── comparar-visual/
│       ├── refinar-codigo/
│       ├── gestionar-estado/
│       └── transferir-conocimiento/
├── commands/              # Comandos slash para acceso rápido
│   ├── analizar-imagen.md
│   ├── generar-tikz.md
│   ├── generar-python.md
│   ├── generar-r.md
│   ├── comparar.md
│   ├── iterar.md
│   ├── exportar.md
│   ├── estado.md
│   ├── auto-iterar.md
│   ├── generar-schoice.md
│   └── generar-cloze.md
└── docs/                  # Documentación del workflow
    ├── 01-EXPLICACION_COMPLETA_GRAFICADOR_EXPERTO.md
    ├── NOMENCLATURA_ARCHIVOS_RMD.md
    ├── INDICE_DOCUMENTACION.md
    ├── WORKFLOW_PASO_A_PASO.md
    └── [otros archivos]
\end{verbatim}

\textbf{Características:}

\begin{itemize}
\tightlist
\item
  ✅ Hooks configurados para recordatorios automáticos
\item
  ✅ Permisos preconfigurados para ejecución sin confirmación
\item
  ✅ Documentación completa del workflow paso a paso
\item
  ✅ Validación automática de estructura ICFES
\item
  ✅ \textbf{Graficador Experto v2.0} con estado persistente y métricas
  cuantitativas
\item
  ✅ Sistema de puntuación 0-100 puntos para comparación visual
\item
  ✅ Transferencia de conocimiento entre lenguajes (TikZ → Python → R)
\item
  ✅ Nomenclatura obligatoria con carpetas estructuradas
\end{itemize}

\begin{center}\rule{0.5\linewidth}{0.5pt}\end{center}

\subsection{\texorpdfstring{🤝 \textbf{Contribución y
Calidad}}{🤝 Contribución y Calidad}}\label{contribuciuxf3n-y-calidad}

Cualquier contribución debe adherirse estrictamente a las metodologías y
protocolos definidos en el proyecto.

\subsubsection{\texorpdfstring{\textbf{Regla de Oro
Anti-Errores}}{Regla de Oro Anti-Errores}}\label{regla-de-oro-anti-errores}

\textbf{``Si no está en los ejemplos funcionales, no improvises.''}
Antes de escribir cualquier código, es \textbf{obligatorio} consultar
los ejercicios en \texttt{/A-Produccion/Ejemplos-Funcionales-Rmd/}.

\subsubsection{\texorpdfstring{\textbf{Criterios de
Calidad}}{Criterios de Calidad}}\label{criterios-de-calidad}

\begin{itemize}
\tightlist
\item
  \textbf{Fidelidad Visual (98\%+)}: Precisión geométrica, cromática, de
  posicionamiento y completitud.
\item
  \textbf{Funcionalidad R-exams (100\%)}: Compatibilidad con
  \texttt{exams2*}, 300+ versiones, y aleatorización completa.
\item
  \textbf{Alineación ICFES}: Metadatos completos, distractores
  pedagógicos y nivel de dificultad apropiado.
\end{itemize}

\begin{center}\rule{0.5\linewidth}{0.5pt}\end{center}

\subsection{\texorpdfstring{📊 \textbf{Información del
Proyecto}}{📊 Información del Proyecto}}\label{informaciuxf3n-del-proyecto}

\begin{itemize}
\tightlist
\item
  \textbf{Autor}: Álvaro Ángel Molina
\item
  \textbf{Institución}: IE Pedacito de Cielo
\item
  \textbf{Propósito}: Generación de ejercicios matemáticos de alta
  calidad para la preparación de la prueba ICFES Saber 11°.
\item
  \textbf{Licencia}: Proyecto Educativo
\item
  \textbf{Última Actualización}: Diciembre 2025
\end{itemize}

\subsection{\texorpdfstring{🆕 \textbf{Novedades
Recientes}}{🆕 Novedades Recientes}}\label{novedades-recientes}

\subsubsection{Graficador Experto v2.0 (Diciembre
2025)}\label{graficador-experto-v2.0-diciembre-2025}

\begin{itemize}
\tightlist
\item
  ✅ \textbf{Estado persistente}: Tracking completo del progreso con
  recuperación ante interrupciones
\item
  ✅ \textbf{Métricas cuantitativas}: Sistema de puntuación 0-100 puntos
  en 6 categorías
\item
  ✅ \textbf{Análisis estructurado}: Formato JSON reutilizable para las
  3 generaciones
\item
  ✅ \textbf{Transferencia de conocimiento}: Lecciones aprendidas
  aplicadas entre lenguajes
\item
  ✅ \textbf{Iteración automática}: Comando \texttt{/auto-iterar} hasta
  umbral de similitud
\item
  ✅ \textbf{Visualización de progreso}: Comando \texttt{/estado} para
  ver avance en tiempo real
\end{itemize}

\subsubsection{Nomenclatura Obligatoria y Carpetas
Estructuradas}\label{nomenclatura-obligatoria-y-carpetas-estructuradas}

\begin{itemize}
\tightlist
\item
  ✅ \textbf{Formato estándar}:
  \texttt{{[}ejercicio{]}\_{[}componente{]}\_{[}competencia{]}\_n{[}nivel{]}\_v{[}version{]}.Rmd}
\item
  ✅ \textbf{Carpetas organizadas}: Cada ejercicio en su propia carpeta
  con todos sus archivos
\item
  ✅ \textbf{Exportación mejorada}: \texttt{/exportar} crea carpeta con
  nomenclatura oficial
\item
  ✅ \textbf{Selección de versión gráfica}: Pregunta obligatoria antes
  de generar .Rmd
\end{itemize}

\subsubsection{Configuración de Tolerancias para Ejercicios
CLOZE}\label{configuraciuxf3n-de-tolerancias-para-ejercicios-cloze}

\begin{itemize}
\tightlist
\item
  ✅ \textbf{Tolerancias apropiadas}: 0 para schoice, ≥1 para valores
  numéricos grandes
\item
  ✅ \textbf{Documentación mejorada}: Guías específicas en comandos y
  skills
\item
  ✅ \textbf{Validación automática}: Verificación de configuración
  correcta
\end{itemize}

Este proyecto representa un sistema integral y robusto que encapsula las
mejores prácticas para la creación de contenido educativo parametrizado
y de alta calidad.

\begin{center}\rule{0.5\linewidth}{0.5pt}\end{center}

\subsection{\texorpdfstring{🧪 \textbf{Ecosistema de Testing
Agresivo}}{🧪 Ecosistema de Testing Agresivo}}\label{ecosistema-de-testing-agresivo}

\subsubsection{🎯 Objetivo: 100\% de
Cobertura}\label{objetivo-100-de-cobertura}

Este proyecto implementa un ecosistema de testing agresivo para
garantizar calidad máxima en la generación automatizada de ejercicios
ICFES.

\begin{longtable}[]{@{}lll@{}}
\toprule\noalign{}
Componente & Cobertura & Estado \\
\midrule\noalign{}
\endhead
\bottomrule\noalign{}
\endlastfoot
\textbf{Validación Matemática} & 100\% & ✅ Completo \\
\textbf{Ortografía Española} & 100\% & ✅ Completo \\
\textbf{Renderizado 4 Formatos} & 100\% & ✅ Completo \\
\textbf{Aleatorización y Diversidad} & 100\% & ✅ Completo \\
\textbf{Flujo B (Graficador)} & 100\% & ✅ Completo \\
\textbf{Tests de Regresión} & 100\% & ✅ Completo \\
\textbf{COBERTURA TOTAL} & \textbf{100\%} & ✅ \textbf{OBJETIVO
ALCANZADO} \\
\end{longtable}

\subsubsection{🧪 Suites de Testing
Implementadas}\label{suites-de-testing-implementadas}

\paragraph{\texorpdfstring{1. \textbf{Validación Matemática}
(\texttt{test\_validacion\_matematica.R})}{1. Validación Matemática (test\_validacion\_matematica.R)}}\label{validaciuxf3n-matemuxe1tica-test_validacion_matematica.r}

\begin{itemize}
\tightlist
\item
  ✅ Detección de errores en chunks R (NaN, Inf, errores de ejecución)
\item
  ✅ Validación de archivos SCHOICE válidos
\item
  ✅ Detección de \texttt{exshuffle\ =\ FALSE} (prohibido)
\item
  ✅ Validación de inconsistencias en CLOZE
\item
  ✅ Verificación de metadatos ICFES completos
\end{itemize}

\paragraph{\texorpdfstring{2. \textbf{Ortografía Española}
(\texttt{test\_ortografia\_espanol.R})}{2. Ortografía Española (test\_ortografia\_espanol.R)}}\label{ortografuxeda-espauxf1ola-test_ortografia_espanol.r}

\begin{itemize}
\tightlist
\item
  ✅ Detección de tildes faltantes
\item
  ✅ Exclusión de metadatos R-exams (ASCII obligatorio)
\item
  ✅ Exclusión de nombres de variables R
\item
  ✅ Correcciones automáticas en texto
\item
  ✅ Preservación de código inline
\end{itemize}

\paragraph{\texorpdfstring{3. \textbf{Renderizado 4 Formatos}
(\texttt{test\_renderizado\_4\_formatos.R})}{3. Renderizado 4 Formatos (test\_renderizado\_4\_formatos.R)}}\label{renderizado-4-formatos-test_renderizado_4_formatos.r}

\begin{itemize}
\tightlist
\item
  ✅ SCHOICE → HTML/PDF/DOCX/NOPS sin errores
\item
  ✅ CLOZE → 4 formatos sin errores
\item
  ✅ Validación cruzada con misma semilla
\item
  ✅ Coherencia de contenido entre formatos
\end{itemize}

\paragraph{\texorpdfstring{4. \textbf{Aleatorización y Diversidad}
(\texttt{test\_aleatorization\_diversity.R})}{4. Aleatorización y Diversidad (test\_aleatorization\_diversity.R)}}\label{aleatorizaciuxf3n-y-diversidad-test_aleatorization_diversity.r}

\begin{itemize}
\tightlist
\item
  ✅ \texttt{exshuffle\ =\ TRUE} genera orden aleatorio
\item
  ✅ Generación de 250+ versiones únicas
\item
  ✅ Cobertura de rangos numéricos esperados
\item
  ✅ Distractores distintos y plausibles
\end{itemize}

\paragraph{\texorpdfstring{5. \textbf{Flujo B Graficador}
(\texttt{test\_flujo\_b\_graficador.R})}{5. Flujo B Graficador (test\_flujo\_b\_graficador.R)}}\label{flujo-b-graficador-test_flujo_b_graficador.r}

\begin{itemize}
\tightlist
\item
  ✅ Estructura de \texttt{workflow\_state.json} correcta
\item
  ✅ Detección obligatoria de gráficos
\item
  ✅ Aprobación secuencial (TikZ → Python → R)
\item
  ✅ Similitud \textgreater= 95\% requerida
\item
  ✅ Verificación de 5 coherencias
\end{itemize}

\paragraph{\texorpdfstring{6. \textbf{Tests de Regresión}
(\texttt{test\_regression\_suite.R})}{6. Tests de Regresión (test\_regression\_suite.R)}}\label{tests-de-regresiuxf3n-test_regression_suite.r}

\begin{itemize}
\tightlist
\item
  ✅ Ejemplos funcionales continúan renderizando
\item
  ✅ Scripts mantienen compatibilidad
\item
  ✅ Hooks mantienen funcionalidad
\item
  ✅ Plantillas mantienen formato
\item
  ✅ Metadatos ICFES siguen estándar
\item
  ✅ Ciclo completo funciona end-to-end
\end{itemize}

\subsubsection{🚀 Ejecución de Tests}\label{ejecuciuxf3n-de-tests}

\paragraph{Ejecutar Suite Completa}\label{ejecutar-suite-completa}

\begin{Shaded}
\begin{Highlighting}[]
\CommentTok{\# Ejecutar todos los tests}
\ExtensionTok{Rscript}\NormalTok{ tests/run\_all\_tests.R}
\end{Highlighting}
\end{Shaded}

\textbf{Salida esperada:}

\begin{verbatim}
========================================
  SUITE DE TESTING COMPLETA
  Repositorio Matemáticas ICFES R-Exams
========================================

Ejecutando: Validación Matemática
--------------------------------------------------
✓ Validación Matemática completado en 2.34 segundos

...

========================================
  REPORTE FINAL
========================================

Suites ejecutadas: 6
✓ Exitosas: 6
✗ Fallidas: 0
Tiempo total: 12.45 segundos

Cobertura de testing: 100.0%
🎉 ¡OBJETIVO DE 100% ALCANZADO!

✅ TODOS LOS TESTS PASARON
\end{verbatim}

\paragraph{Ejecutar Suite Individual}\label{ejecutar-suite-individual}

\begin{Shaded}
\begin{Highlighting}[]
\CommentTok{\# Solo validación matemática}
\ExtensionTok{Rscript} \AttributeTok{{-}e} \StringTok{"library(testthat); test\_file(\textquotesingle{}tests/testthat/test\_validacion\_matematica.R\textquotesingle{})"}

\CommentTok{\# Solo ortografía}
\ExtensionTok{Rscript} \AttributeTok{{-}e} \StringTok{"library(testthat); test\_file(\textquotesingle{}tests/testthat/test\_ortografia\_espanol.R\textquotesingle{})"}

\CommentTok{\# Solo renderizado}
\ExtensionTok{Rscript} \AttributeTok{{-}e} \StringTok{"library(testthat); test\_file(\textquotesingle{}tests/testthat/test\_renderizado\_4\_formatos.R\textquotesingle{})"}
\end{Highlighting}
\end{Shaded}

\subsubsection{🔄 Integración Continua
(CI/CD)}\label{integraciuxf3n-continua-cicd}

El repositorio incluye configuración de GitHub Actions
(\texttt{.github/workflows/ci-testing.yml}) que ejecuta automáticamente:

\textbf{Triggers:} - ✅ Cada push a \texttt{main} o \texttt{develop} -
✅ Cada pull request - ✅ Diariamente a las 02:00 UTC

\textbf{Jobs Paralelos:} 1. Tests de Validación Matemática 2. Tests de
Ortografía 3. Tests de Renderizado 4 Formatos 4. Tests de Aleatorización
y Diversidad 5. Tests de Flujo B (Graficador) 6. Tests de Regresión 7.
Reporte de Cobertura

\textbf{Política:} Tolerancia cero a regresiones. Si algún test falla,
el pipeline completo falla.

\subsubsection{🛡️ Política de Testing}\label{poluxedtica-de-testing}

\paragraph{Reglas Obligatorias}\label{reglas-obligatorias}

\begin{itemize}
\tightlist
\item
  ❌ \textbf{PROHIBIDO} hacer push a \texttt{main} si algún test falla
\item
  ❌ \textbf{PROHIBIDO} usar \texttt{git\ commit\ -\/-no-verify} para
  evadir hooks
\item
  ❌ \textbf{PROHIBIDO} comentar tests que fallan (arreglar el código,
  no el test)
\item
  ❌ \textbf{PROHIBIDO} reducir cobertura por debajo del 100\%
\end{itemize}

\paragraph{Proceso de Contribución}\label{proceso-de-contribuciuxf3n}

\begin{enumerate}
\def\labelenumi{\arabic{enumi}.}
\tightlist
\item
  \textbf{Antes de cada commit:} Hook pre-commit valida ortografía
  automáticamente
\item
  \textbf{Después de cada \texttt{exams2*()}:} Hook post-exams2 valida
  matemática + genera preview PNG
\item
  \textbf{Antes de cada push:} Ejecutar
  \texttt{Rscript\ tests/run\_all\_tests.R} localmente
\item
  \textbf{Antes de cada merge:} CI/CD ejecuta suite completa
  automáticamente
\end{enumerate}

\subsubsection{📊 Métricas de Calidad}\label{muxe9tricas-de-calidad}

\begin{longtable}[]{@{}llll@{}}
\toprule\noalign{}
Componente & Tests & Tiempo Promedio & Crítico \\
\midrule\noalign{}
\endhead
\bottomrule\noalign{}
\endlastfoot
Validación Matemática & 5 tests & \textasciitilde2.5s & ✅ Sí \\
Ortografía Española & 5 tests & \textasciitilde1.5s & ✅ Sí \\
Renderizado 4 Formatos & 6 tests & \textasciitilde8.0s & ✅ Sí \\
Aleatorización & 4 tests & \textasciitilde5.0s & ✅ Sí \\
Flujo B & 6 tests & \textasciitilde1.0s & ✅ Sí \\
Regresión & 7 tests & \textasciitilde10.0s & ✅ Sí \\
\textbf{TOTAL} & \textbf{33+ tests} & \textbf{\textasciitilde28.0s} &
- \\
\end{longtable}

\subsubsection{📚 Documentación
Completa}\label{documentaciuxf3n-completa}

Ver documentación detallada del ecosistema en: -
\textbf{\texttt{.claude/docs/ECOSISTEMA\_TESTING.md}} - Guía completa de
testing - \textbf{\texttt{tests/testthat/}} - Suites de tests
individuales - \textbf{\texttt{tests/run\_all\_tests.R}} - Script
ejecutor principal - \textbf{\texttt{.github/workflows/ci-testing.yml}}
- Configuración CI/CD

\begin{center}\rule{0.5\linewidth}{0.5pt}\end{center}

\subsection{\texorpdfstring{🔗 \textbf{Integración Graficador Experto
v2.0 con
R-exams}}{🔗 Integración Graficador Experto v2.0 con R-exams}}\label{integraciuxf3n-graficador-experto-v2.0-con-r-exams}

El sistema incluye un flujo completo e integrado para transformar
gráficos matemáticos en ejercicios R-exams:

\subsubsection{Flujo Completo: Imagen → Gráfico →
Ejercicio}\label{flujo-completo-imagen-gruxe1fico-ejercicio}

\begin{verbatim}
1. ANÁLISIS VISUAL
   /analizar-imagen grafico.png
   ↓
   Genera: analisis_inicial.json + workflow_state.json

2. GENERACIÓN MULTI-LENGUAJE
   /generar-tikz → Validado (98%+)
   /generar-python → Validado (95%+)
   /generar-r → Validado (96%+)
   ↓
   Genera: output_tikz.tex + output_python.py + output_r.R

3. COMPARACIÓN CON MÉTRICAS OBJETIVAS
   /comparar → Puntuación 0-100 en 6 categorías
   /auto-iterar r 95 10 → Iteración automática
   ↓
   Resultado: 95+ puntos de similitud garantizados

4. EXPORTACIÓN CON NOMENCLATURA
   /exportar
   ↓
   Pregunta: ¿Qué versión usar? (TikZ/Python/R)
   ↓
   Crea: [ejercicio]_[componente]_[competencia]_n[nivel]_v[version]/
         ├── Códigos de las 3 versiones
         ├── Imágenes generadas
         ├── Análisis y estado persistente
         └── Reporte consolidado

5. GENERACIÓN DE EJERCICIO R-EXAMS
   /generar-schoice (o /generar-cloze)
   ↓
   Usa la versión gráfica seleccionada
   ↓
   Genera: [nombre_completo].Rmd dentro de la carpeta
\end{verbatim}

\subsubsection{Ventajas de la
Integración}\label{ventajas-de-la-integraciuxf3n}

\begin{longtable}[]{@{}
  >{\raggedright\arraybackslash}p{(\linewidth - 2\tabcolsep) * \real{0.4091}}
  >{\raggedright\arraybackslash}p{(\linewidth - 2\tabcolsep) * \real{0.5909}}@{}}
\toprule\noalign{}
\begin{minipage}[b]{\linewidth}\raggedright
Ventaja
\end{minipage} & \begin{minipage}[b]{\linewidth}\raggedright
Descripción
\end{minipage} \\
\midrule\noalign{}
\endhead
\bottomrule\noalign{}
\endlastfoot
\textbf{Trazabilidad} & Cada ejercicio tiene historial completo de cómo
se generó su gráfico \\
\textbf{Flexibilidad} & 3 versiones disponibles según necesidad
(vectorial/Python/R) \\
\textbf{Calidad} & Métricas objetivas 0-100 puntos antes de usar en
.Rmd \\
\textbf{Organización} & Carpetas nomencladas con todos los archivos
relacionados \\
\textbf{Reproducibilidad} & Estado persistente permite regenerar o
modificar \\
\end{longtable}

\subsubsection{Recomendaciones por Tipo de
Gráfico}\label{recomendaciones-por-tipo-de-gruxe1fico}

\begin{longtable}[]{@{}
  >{\raggedright\arraybackslash}p{(\linewidth - 4\tabcolsep) * \real{0.3778}}
  >{\raggedright\arraybackslash}p{(\linewidth - 4\tabcolsep) * \real{0.4667}}
  >{\raggedright\arraybackslash}p{(\linewidth - 4\tabcolsep) * \real{0.1556}}@{}}
\toprule\noalign{}
\begin{minipage}[b]{\linewidth}\raggedright
Tipo de Gráfico
\end{minipage} & \begin{minipage}[b]{\linewidth}\raggedright
Versión Recomendada
\end{minipage} & \begin{minipage}[b]{\linewidth}\raggedright
Razón
\end{minipage} \\
\midrule\noalign{}
\endhead
\bottomrule\noalign{}
\endlastfoot
Geometría precisa & TikZ & Vectorial, máxima precisión \\
Estadística básica & R/ggplot2 & Nativo R-exams, fácil mantener \\
Visualización compleja & Python/matplotlib & Flexibilidad, numpy
integrado \\
Funciones matemáticas & TikZ o R & Calidad y precisión \\
Gráficos de barras/líneas & R/ggplot2 & Sintaxis simple, temas
profesionales \\
\end{longtable}

\begin{center}\rule{0.5\linewidth}{0.5pt}\end{center}

\end{document}
